\section{Reproducibility and metric definitions}
\label{app:reproducibility}

The analysis artifact contains execution-level records for the
model-written policies and the exploratory runtime-record comparison,
alongside final aggregate summaries for the hybrid policies.
Table entries and numerical claims are generated from those
inputs. The aggregate summaries contain solved-rate margins and matched
discordant outcomes, sufficient to reconstruct the reported success
contrasts and exact McNemar tests. They do not contain complete
trajectories, per-call prompt measurements, or joint task-level outcomes
for all configurations.

Resource ratios are computed separately for each condition against its
own full-history reference, restricting to shared successes. For a
resource $X$, the reported statistic is
\[
 \operatorname{median}_{(i,r):\,Y_{iar}=Y_{i0r}=1}
 \frac{X_{iar}}{X_{i0r}}.
\]
A zero reference denominator would make an individual ratio undefined.
The analysis does not replace such a denominator with an arbitrary
constant. Input and cached tokens are separate recorded counters.

The efficiency measures use all retained outcomes:
\[
 \text{Solved/\$}=\frac{\sum Y}{\sum D},
 \qquad
 \text{Solved/Mtok}=
 \frac{10^{6}\sum Y}{\sum(U+K+O)}.
\]
The latter measures logical token volume under the recorded counter
convention, not equivalence of computational work across cached and
uncached tokens. Hybrid cost efficiencies are approximate because the
available aggregate costs are rounded. Unavailable token totals and
prompt peaks are not reconstructed from marginal medians.

\section{Sensitivity of exploratory component costs}
\label{app:composition}

The exploratory comparison permits a check of outcome-conditioned
selection. A separate shared-success subset for each configuration changes
the tasks contributing to its cost ratio. To compare the configurations
on the same tasks, we additionally restrict all ratios to the intersection
of successes across every included configuration and the control.
Table~\ref{tab:composition-sensitivity} reports both calculations.

\begin{table}[htbp]
\centering\small
\caption{Exploratory resource sensitivity. Variable-subset cost ratios use
each configuration's shared successes with full history. Fixed-subset
ratios use the same all-configuration success intersection. Entries are
medians of task-wise ratios to full history, not ratios of marginal
medians.}
\label{tab:composition-sensitivity}
\input{tables/new-composition-sensitivity}
\end{table}

The two richest representations have variable-subset cost ratios of
\ExploreRichVariableCostRatio{} and
\ExploreFailuresVariableCostRatio{}, but fixed-subset ratios of
\ExploreRichFixedCostRatio{} and \ExploreFailuresFixedCostRatio{}.
The large reduction relative to file references remains. A cost advantage
over full history does not. The fixed subset is still selected on outcomes,
so this sensitivity check describes successful-task resource use rather
than an unconditional causal effect.

\section{Complete efficiency breakdown}
\label{app:efficiency}

\begin{table}[htbp]
\centering\small
\caption{Cost and efficiency over retained executions, including failures.
Successes per dollar use aggregate cost and observed successes, not ratios
of medians. Token efficiency is reported where cumulative totals are
available. Final hybrid efficiencies use rounded reported cost totals.}
\label{tab:economics}
\input{tables/new-economics}
\end{table}

Table~\ref{tab:economics} supplements the cost-per-success comparisons in
the main text. Median execution cost and aggregate cost per success answer
different questions: the latter also includes the cost of failures and the
shape of the cost distribution. Neither estimates the cost of repeated
attempts until success.

\section{Success by reference trajectory length}
\label{app:horizon}

Table~\ref{tab:horizon} separates tasks using the full-history control's
median recorded trajectory length. All state policies use the same
partition. The larger state-only deficit appears in the longer-reference
subset, while recent-history retention improves the observed solved rate
in both subsets. Because the partition is post-hoc and trajectory length
also reflects task characteristics, this pattern does not identify a
causal horizon effect or a context-window crossover.

\begin{table}[htbp]
\centering\small
\caption{Solved rate by the full-history control's median recorded turns
$T$. Differences and exact McNemar tests use full history within the same
subset. The grouping is descriptive and held fixed across conditions.}
\label{tab:horizon}
\input{tables/new-horizon}
\end{table}
